%! Author = chtompki
%! Date = 1/14/26
% Example LaTeX document for GP111 - note % sign indicates a comment
\documentclass[12pt]{amsart}
% Default margins are too wide all the way around. I reset them here
\setlength{\hoffset}{-.75in}
\setlength{\textwidth}{6.5in}
\setlength{\voffset}{-.5in}
\setlength{\textheight}{9.0in}
\usepackage{amsmath}
\usepackage{amssymb}
\usepackage{amsthm}
\usepackage{pgfplots}
\usepackage{enumerate}
\usepackage{tikz}
\usetikzlibrary{shadows,arrows,arrows.meta,shapes,positioning,calc}% Required for the Circle and Triangle arrow tips
\usepackage{setspace}
\usepackage{siunitx}
\usepackage{enumitem}

\theoremstyle{conjecture}
\newtheorem{conjecture}{Conjecture}

\theoremstyle{definition}
\newtheorem{definition}{Definition}

\theoremstyle{question}
\newtheorem{question}{Question}

\theoremstyle{theorem}
\newtheorem{theorem}{Theorem}

\newenvironment{solution}
{\begin{proof}[Solution]}
{\end{proof}}

\title{Graph Theory:\\Homework 7\\ Book Section 3.1 Part III and 4.1.}
\author{Rob Tompkins}
\date{\today}


\begin{document}

    \maketitle
    \noindent\textbf{1.} Let $M$ be a maximum matching and $L$ a minimal edge cover in a graph with no
    isolated vertices.
    Probe the following statements below:
    \begin{itemize}
        \item[(a)] $M$ is a maximal matching if and only if $M$is contained in a minimum edge cover.
        \item[(b)] $L$ is a minimum edge cover if and only if $L$ contains a maximal matching.
    \end{itemize}

    \begin{proof}

    Write $n=|V(G)|$, let $\nu$ be the maximum matching size, and let
    $\rho$ be the minimum edge cover size. We first show that
    $\rho=n-\nu$.

    Every minimal edge cover $F$ induces a disjoint union of stars, each
    with at least one edge. Indeed, an edge whose two endpoints both have
    degree at least two in $(V(G),F)$ could be deleted without uncovering
    a vertex. Thus every edge has a degree-one endpoint, which forces
    each component to be a star. If there are $c$ such components, then
    $|F|=n-c$, and choosing one edge from each star gives a matching of
    size $c$. Consequently $|F|\ge n-\nu$. Applying this to a minimum
    edge cover gives $\rho\ge n-\nu$.

    Conversely, start with a maximum matching $M$. Its unmatched vertices
    form an independent set, since an edge between two of them would
    enlarge the matching. Choose an incident edge for each unmatched
    vertex; this is possible because there are no isolated vertices.
    These edges are distinct and join unmatched vertices to matched
    vertices. Together with $M$ they form an edge cover of size
    $\nu+(n-2\nu)=n-\nu$. Hence $\rho=n-\nu$.

    \smallskip
    \noindent\textbf{(a)} Under the opening assumption that $M$ is a
    \emph{maximum} matching, both assertions are true: $M$ is maximal,
    and the preceding construction extends it to a minimum edge cover.
    Thus the stated equivalence holds under that assumption.
    If the intended statement instead quantifies over arbitrary matchings,
    ``maximal'' must be replaced by ``maximum.'' For example, on the path
    $1,2,3,4$, the matching $\{23\}$ is maximal, but the unique minimum
    edge cover is $\{12,34\}$.

    For completeness, a matching contained in a minimum edge cover need
    not itself be maximum (it could be empty). The usual unrestricted
    equivalence requires the matching to be \emph{maximal}: a maximal
    matching is maximum if and only if it is contained in a minimum edge
    cover. To see the reverse direction, a minimum edge cover has exactly
    $\nu$ star components. A maximal matching contained in it must choose
    one edge from every star, and therefore has size $\nu$.

    \smallskip
    \noindent\textbf{(b)} As written, this statement is false. On the path
    $1,2,3,4,5,6$, let
    \[
       L=\{12,23,45,56\}.
    \]
    This is a minimal edge cover, since deleting any of its edges uncovers
    a leaf of one of its two star components. It contains the maximal
    matching $\{12,45\}$: the only unmatched vertices are $3$ and $6$,
    which are not adjacent. However, $L$ is not minimum, since
    $\{12,34,56\}$ is an edge cover with only three edges.

    The corrected assertion replaces ``maximal'' by ``maximum.'' If $L$
    is minimum, its star decomposition has $n-|L|=\nu$ components, so
    selecting one edge from each gives a maximum matching contained in
    $L$. Conversely, if the minimal edge cover $L$ contains a maximum
    matching, each of its $c$ star components contributes at most one
    matching edge. Hence $\nu\le c\le\nu$, so
    $|L|=n-c=n-\nu=\rho$, and $L$ is minimum.
    \end{proof}

    \noindent\textbf{2.} Let $G$ be a simple graph in which the sum of degrees of any $k$ vertices
    is less than $n-k$.
    Prove that every maximal independent set in $G$ has more than $k$ vertices.

    \begin{proof}

    Let $I$ be a maximal independent set, and suppose $|I|\le k$.
    Choose a set $S$ of $k$ vertices containing $I$. (As usual the
    hypothesis is understood with $0\le k\le n$.) Maximality implies
    that every vertex outside $I$, and hence every vertex outside $S$,
    has a neighbor in $I$. Thus at least $n-k$ edges join $I$ to
    $V(G)\setminus S$, giving
    \[
       n-k\le\sum_{v\in I}d(v)\le\sum_{v\in S}d(v)<n-k,
    \]
    a contradiction. Therefore $|I|>k$.
    \end{proof}

    \noindent\textbf{3.} An edge $e$ of a graph $G$ is $\alpha$-critical is $\alpha(G - e) > \alpha(G)$.
    Suppose that $xy$ and $xz$ are $\alpha$-critical edges in $G$.
    Prove that $G$ has an induced subgraph that is an odd cycle containing $xy$ and $xz$.
    (Hint: Let $Y$, $Z$ be maximum independent sets in $G-xy$ and $G-xz$ respectively.
    Let $H = G[Y\Delta Z]$.
    Prove that every component of $H$ has the same number of vertices from $Y$ to $Z$.
    Use this to prove $y$ and $z$ belong to the same component of $H$.)

    \begin{proof}

    Put $a=\alpha(G)$. Deleting one edge increases the independence
    number by at most one: removing one endpoint of that edge from an
    independent set in the edge-deleted graph leaves an independent set
    in $G$. Therefore the sets in the hint satisfy
    $|Y|=|Z|=a+1$. Moreover $x,y\in Y$ and $x,z\in Z$; otherwise
    one of these sets would already be independent in $G$. Since $xz$
    is present in $G-xy$ and $xy$ is present in $G-xz$, we have
    $z\notin Y$ and $y\notin Z$.

    Set $A=Y\setminus Z$, $B=Z\setminus Y$, and $H=G[A\cup B]$.
    The sets $Y\setminus\{x\}$ and $Z\setminus\{x\}$ are independent
    in $G$ and have size $a$, so both are maximum. In particular $H$
    is bipartite with parts $A,B$.

    Let $C$ be any component of $H$. Exchanging $A\cap V(C)$ for
    $B\cap V(C)$ in $Y\setminus\{x\}$ preserves independence:
    there are no edges to another component of $H$, and vertices of
    $(Y\cap Z)\setminus\{x\}$ have no neighbors in either part.
    Maximality of the cardinality $a$ implies
    $|B\cap V(C)|\le |A\cap V(C)|$. Exchanging in the other direction
    using $Z\setminus\{x\}$ gives the reverse inequality. Hence each
    component has equally many vertices in $A$ and $B$.

    Suppose $y$ and $z$ lie in different components, and let $C$ contain
    $y$. Then
    \[
       Y'=(Y\setminus(A\cap V(C)))\cup(B\cap V(C))
    \]
    has size $a+1$. The same exchange argument shows that its vertices
    other than $x$ are independent. The only neighbor of $x$ in $Y$
    is $y$, which was removed, and the only neighbor of $x$ in $Z$ is
    $z$, which was not added. Thus $Y'$ is independent in $G$, a
    contradiction. Therefore $y$ and $z$ belong to the same component.

    Choose a shortest $y$--$z$ path $P$ in $H$. Because $y\in A$ and
    $z\in B$, its length is odd. A shortest path has no chord, and $x$
    has no neighbors in $V(H)\setminus\{y,z\}$. Consequently $P$
    together with the edges $xy,xz$ is an induced cycle of odd length,
    containing both specified edges.
    \end{proof}

    \noindent\textbf{4.} Let $G$ be a connected graph with at least three vertices.
    Form $G'$ and $G$ by adding an edge with endpoints $x,y$ whenever $d_G(x,y) = 2$.
    Prove that $G'$ is 2-connected.

    \begin{proof}

    Interpret the construction as adding the specified edges to $G$.
    Fix any vertex $v$. For any two distinct vertices $a,b\ne v$,
    choose a simple $a$--$b$ path in $G$. If it avoids $v$, it remains
    a path in $G'-v$. Otherwise let $u,v,w$ be the consecutive vertices
    where it passes through $v$. The distinct vertices $u,w$ are either
    adjacent in $G$ or have distance two in $G$, so $uw\in E(G')$.
    Replacing $u,v,w$ by $u,w$ gives an $a$--$b$ path in $G'-v$.
    Thus $G'-v$ is connected for every $v$. Since $G'$ is connected
    and has at least three vertices, it is 2-connected.
    \end{proof}

    \noindent\textbf{5.} Prove that $\kappa(G) = \delta(G)$ if $G$ is simple and $\delta(G) \ge n(G) - 2$.
    Prove that this is the best possible for each $n\ge 4$ by constructing a simple $n$-vertex graph with
    minimum degree $n-3$ and connectivity less than $n-3$.

    \begin{proof}

    Write $n=n(G)$, using the standard convention
    $\kappa(K_n)=n-1$. If $G$ is complete, the desired equality is
    immediate. Otherwise $\delta(G)=n-2$, and each vertex has at most
    one nonneighbor among the other vertices.

    Suppose $S$ is a vertex set with $|S|<n-2$. At least three vertices
    remain. If $G-S$ were disconnected, choose vertices $u,v$ in two
    distinct components, and a third remaining vertex $w$. If $w$ is
    in the component of $u$, then $v$ has the two nonneighbors $u,w$;
    otherwise $u$ has the two nonneighbors $v,w$. Both are impossible.
    Hence deleting fewer than $n-2$ vertices cannot disconnect $G$.
    On the other hand, since $G$ is not complete, choose nonadjacent
    vertices $p,q$. Deleting the other $n-2$ vertices leaves two isolated
    vertices. Thus $\kappa(G)=n-2=\delta(G)$.

    For sharpness, for each $n\ge4$ take two disjoint copies of $K_2$
    and a clique $S$ on $n-4$ vertices, and join every vertex of $S$
    to all four vertices in the two copies of $K_2$. Each of those four
    vertices has degree $1+(n-4)=n-3$, while each vertex in $S$ has
    degree $n-1$. Therefore $\delta(G)=n-3$. Deleting $S$ leaves
    $K_2\mathbin{\dot\cup}K_2$, so
    \[
       \kappa(G)\le n-4<n-3=\delta(G).
    \]
    When $n=4$, $S$ is empty and the graph is already disconnected,
    with connectivity zero. Thus the construction covers every $n\ge4$.
    \end{proof}


\end{document}